\chapter*{前言}
\addcontentsline{toc}{chapter}{前言}
\markboth{前言}{前言}

中國目前的文字學研究，
絕大多數情況下並不能稱爲真正的文字學研究，
而更適合稱爲「漢字學」（Sinography），
甚至只能說是「古漢字學」（Chinese palaeography）。
這並不算什麼大的問題，
只是佔用了一個真正可能發展爲對「文字」進行研究的學科的專名。
我不好說這對中國的文字研究產生了什麼影響，
至少我所瞭解的在中國專門研究世界文字的人寥寥無幾。
因此，
本講義既不能太前沿，
也不能繼續走國內傳統的小學路子——我們會嘗試引入極端而理想的理論模型，
將語言結構與數學結構對應起來；但是在實踐上，
我們又必須強調這樣的理論模型是完全不符合實際的，
於是又需要以保守的態度對待之。


所謂文字學（graphemics），
應該是研究全體文字的學科，
而不僅限於漢字。
我們國內的研究漢字及古漢字的學科，
只能夠視爲文字學的分支。
文字學作爲一個學科，
所面對的是世界上所存在的所有（自然的）文字，
所要完成的任務首先是瞭解文字的特性，
以及我們如何認識其內部的結構，
而最重要的認識工具即是分類。
分類有兩個層面的，
即從索緒爾（Saussure）開始區分的「共時」（synchronic）和「歷時」（diachronic）兩條路，
前者實際上是我們所稱爲「文字類型學」的部分，
後者則是「文字史」的主要討論內容。


文字（script）或者書寫系統（writing system），
本質只是一種符號（sign），
而且是一種視覺符號。
但是，
自然的文字是沒有直接映射到意義（meaning）層的能力的，
必須以語言（language）作爲一層中介。
因此，
文字是作爲一種語言符號的符號來映射到意義的。
理想狀況下，
文字（$\mathcal S$）、 語言（$\mathcal L$）、意義（$\mathcal M$）構成兩層映射：
\[\varphi:\ \mathcal S\rightarrow \mathcal L,\ \psi:\ \mathcal L\rightarrow \mathcal M\]則文字表達意義是通過$\chi=\psi\circ\varphi:\mathcal S\rightarrow \mathcal M$作用的。
如果只討論一個相當理想的模型（當然我們知道，
現實完全不是這樣），
上述映射都應該是同構（isomorphism）——即其都按照相應的組合規則（例如文字的線性組合、語言的語法規則、意義的邏輯關係）形成結構上的一一對應。
如果意義是底層，
語言是第一層，
則文字是從屬於語言的第二層——草率地說，
文字是語言的附庸（語言和意義的關係就沒那麼簡單了，
但那是哲學家們該思考的問題）。


文字類型學（typology of scripts），
即在共時的層面上對文字進行分類。
分類的依據是共時的，
而完全不需要考慮其歷史演變。
我們對類型學的瞭解遠遠不夠，
常常受到漢字構型學（formation of Chinese characters）的強烈影響，
很多著名學者對類型學都有一定的誤解。
索緒爾提出的表音文字（phonogram）、表意文字（ideogram）二分的結構是很有洞見的，
至少在處理大頭上不會犯錯。
比較前沿的類型學研究關注更細節的東西，
涉及正字法的「深度」（depth）以及實際運用中的不同「層」（layer）——這些是精細化研究的結果，
但我們的理論還沒有建構到那一層深度，
應當先着眼於傳統的類型學方法——即以文字單位與語言單位的映射特徵爲依據。
按結構主義（structuralism）語言學的觀點，
共時層是先於歷時層的，
所以我們講義的結構還是以類型學的視角展開。


文字史（history of writing systems），
即研究文字演進的歷史。
這一點跟通常的歷史沒什麼兩樣，
只不過這裏我們或許更強調一個書寫系統到另一個書寫系統的演進，
但也可以涉及一些系統內部更微小的變化——前者就像是從腓尼基文演化成希臘文，
後者就像是漢字的簡化。
和語系（language family）或者生物分類學（taxonomy）一樣，
這種文字史的研究多半仰賴系統發生學（phylogenetics）方法——不過不像生物學還可以按生殖隔離（reproductive isolation）來判斷，
文字的劃分跟語言一樣很難找到一個完全客觀的標準：比如我們可以認爲漢語的漢字和越南的喃字是同一種文字（畢竟喃字的部件和構型我們實在是太熟悉了），
但也可以認爲他們完全是兩種文字（畢竟未經訓練的漢語漢字使用者肯定是一個喃字都看不懂的）。
我們想把這個問題按下不表——這不影響我們的討論。
值得高興的是，
不像文字構型學那樣，
我們好歹還能看見周有光先生的大作《世界文字發展史》，
使得我國在文字史方面不至於太空白。
文字史當然是很重要的，
但本講義不可能全部覆蓋。


最後應該大概闡述一下本講義的組織綱要。
我們的目的是對文字學有一個入門級別的介紹，
選擇的觀點不會太新，
但至少經典，
不會犯太多錯。
核心是以文字類型學爲大框架，
依次介紹全音位文字、輔音音位文字、元音附標文字、音節文字和語素文字。
按索緒爾的兩分法，
前四者屬於表音文字，
最後一個屬於表意文字。
對於每一類文字，
我們都會首先簡要介紹其特徵，
列舉哪些文字是這一類的，
然後選擇一些具體的文字案例進行比較詳細的介紹。


選取哪些文字案例來詳細介紹，
除了考慮其在歷史上的重要性或者文字類型上的代表性之外，
最好還不會太難，
也不至於與讀者太遙遠——這就讓我們不會去選擇那些雖然很重要但是已經滅絕太久的文字來過多介紹了：比如我們更傾向於介紹希臘文、拉丁文，
而不是去介紹原始西奈文、原始迦南文；我們更傾向於介紹天城文，
而不是去介紹婆羅米文——當然我們會在文字史的部分提到他們。


文字案例的詳細介紹，
主要是包括其基本的特徵，
字母表，
以及相關的文字史。
我們會嘗試讓大家去記住一些文字系統，
表音文字中重點講到的案例當然最好能夠全部記住。
文字作爲語言的附庸，
就意味着儘管我不想去談任何與語法相關的東西，
但還是不得不嘗試引入語言來讓大家能夠「讀」出來這些文字。
選取經典語言是必要的，
就算僅僅是方便大家記憶這些符號而已。
這樣的語言一般不會是現代語言，
因爲現代語言的正字法（即拼寫的方式）很可能已經沒那麼「言文一致」了（想想英語就知道了）——爲了讓大家記了字母就能讀出單詞甚至文章，
就必須盡可能選擇言文對應比較好的語言，
而這樣的語言往往都是某種文字最開始使用時候的樣貌。
但是，
提到現代語言又是實操上不可避免的，
所以我們也許可以介紹一下（還能讓大家看看像藏文這種語言發生巨大變化而正字法幾乎不變的案例）。